% Cote 119: offprints, reports, courses, CV and letters (typed and printed part): archive cote 119, site page 164 (PDF page 165).
% typescript. Draft compared: none.
% AI-transcribed by Claude Opus 5.5 (claude-opus-5-5, adaptive thinking, effort high), Message Batches API, batch r119_full (msgbatch_01DwZNm5Y7YsKBWyAN4sFs8w), 2026-11-13; not yet checked by eye.
\scanpage{119/164}
\mspage{- 3 -}

Les fonctions indéfinimen\unsure{t} dérivables sont donc assez peu rigides
pour approximer n'importe qu'elle fonction continue (ici à support compact).
On pourrait obtenir également un résultat d'approximation sur un compact par
des fonctions analytiques (et même polynômes) mais le principe du prolongement
analytique nous interdit d'espérer un résultat sur les supports des fonctions
approximantes.
De plus si l'on quitte le domaine réel ce dernier résultat devient lui-même
faux et il ne sera plus question d'approximer une fonction continue de $\mathbb{C}$
dans $\mathbb{C}$ par des fonctions analytiques.

2 . \uline{Intégrales curvilignes}

\uline{Définition 2.1}. On appelle chemin différentiable une application $t \to \gamma(t)$
du segment $[a,b]$ dans le plan $\mathbb{R}^2$ telle que les coordonnées $x(t)$ et $y(t)$
du point $\gamma(t)$ soient des fonctions continûment différentiables. On supposera
$a < b$ . L'origine de $\gamma$ est le point $\gamma(a)$, l'extrémité le point $\gamma(b)$.
Si $U$ est un ouvert du plan, on dit que $\gamma$ est un chemin différentiable de
l'ouvert $U$ si la fonction $\gamma$ prend ses valeurs dans $U$.

\uline{Définition 2.2}. Une forme différentiable dans un ouvert $U$ du plan est une
expression $\omega = P\,dx + Q\,dy$ dont les coefficients $P$ et $Q$ sont des fonctions
continues dans $D$. (à valeurs dans $\mathbb{R}$ ou dans $\mathbb{C}$). Si $\gamma$ est un chemin
différentiable de $U$ et $\omega$ une forme différentielle dans $U$, on définit
l'intégrale $\int_\gamma \omega$ par la famille $\int_a^b \gamma^*(\omega)$

où $\gamma^*(\omega)$ désigne la forme différentielle $f(t)\,dt$ définie par
$f(t) = P(x(t), y(t))\, x'(t) + Q(x(t), y(t))\, y'(t)$. Autrement dit $\gamma^*(\omega)$ est
la forme différentielle déduite de $\omega$ par le changement de variable $x = x(t\unsure{))|)}$
$y = y(t)$. Si on change de paramètre : $t = t(u)$ [c'est-à-dire, par définition,
$u \to t(u)$ est une fonction \uline{continûment dérivable} pour $a_1 \leq u \leq b_1$ dont la

\mspage{164}
